\section{おわりに}\label{sec6}
\subsection{まとめ}\label{sec6.1}
本研究では，物体の状態・ID・関係を用いる家庭内タスクの逐次計画を対象に，知識の抽出・自然言語化，事例検索，LLMによる前提条件判定と1回の再生成，終了・採点を追跡できる実装を示した．その評価を基本性能，構成要素比較，ケース分析の流れで行った．前提条件検証・再計画の一般原理そのものを新規な技術とは位置付けず，LLM判定と記号判定を比較しながら，具体的な実装の働きと制約を検討した．

データは状態変化・配置の計618インスタンスで，テスト415件と事例用203件からなる．定型指示の下で同名物体を識別し，初期状態に応じて行動を変更する設定であり，既存ベンチマークを補完する診断対象とした．シーン・操作構造は事例とテストで共有し，未知家屋への汎化や一般的な人間難易度を測るものではない．

基本性能評価の18条件では，GPT-4o・Multi・検証ありの成功数は状態変化279/312（89.4\%），配置84/103（81.6\%）だった．成功・失敗別の列長を比較することで，途中終了による短い行動列と，成功時の手順の違いを区別した．ただし，条件間の入力差と実行記録の制約があるため，成功率の差を前提条件検証だけの因果効果とは解釈しない．

構成要素比較は4タスク・2モデル・10条件・3反復の240エピソードである．LLM検証あり／なしの成功数はminiで9/12対8/12，4oで9/12対9/12で，各分子に初期充足タスクの3反復を含む．記号検証，一般的再考，失敗後修復，検索変更との比較も一律の優位性を示さなかった．実行ログでは，条件判定，未達なのにEndを返す判断，JSON条件を満たしても生じるVH失敗を分離した．棄却候補の再実行には状態一致の制約があり，違反説明の著者確認もAI支援後の記入で，説明品質を独立に検証したものではない．これらを踏まえ，本研究の知見は具体的な実装の失敗過程を示す診断的証拠に限定する．


\subsection{今後の課題}\label{sec6.2}
まず，接頭列方式と人物のいないシーンへの復元方式を用いても状態が不一致だった4件の再現と，違反説明に対する評価基準・判断根拠の妥当性の検証が課題である．本評価は照合通過4候補の結果と未再現範囲の開示に限定しており，この4件のVH成功，事後照合を伴う再実行のVH成功6/8，著者の「正しい」1件・「誤り」7件という記入のいずれも，母集団の誤棄却率や違反説明の確定正解率として代用しない．

次に，ベースラインの設定とデータ構築手順を含む再現性を確保し，実行条件を統制した全415タスクの反復比較，共通候補での判定比較，入力・出力・総要求数を考慮した計算量対照を行う必要がある．状態情報を加える検索や再生成後の再検証も検討対象だが，本評価の仕様には含まれず，独立した条件として評価する必要がある．

さらに，定型指示と共有シーンを越える課題，より複雑な物体相互作用や未知環境への拡張が必要である．実機への適用には，センサーからの認識誤差，低レベル制御，ゴール以外の安全制約も含めた別の評価を要する．本研究のシミュレータ上の成功率や条件辞書の充足を，そのまま実世界の信頼性・安全性の保証とはしない．
